% ==============================================================================
% CONFIGURACIÓN GLOBAL Y PAQUETES (UNAL - DIEEC)
% ==============================================================================

% --- PAQUETES BÁSICOS Y IDIOMA ---
\usepackage[utf8]{inputenc}
\usepackage[spanish,es-nodecimaldot]{babel}
\usepackage[margin=1.5cm, top=2.2cm, bottom=2.2cm, headheight=15pt]{geometry}

% --- PAQUETES DE FORMATO Y ESTILO ---
\usepackage{xcolor}
\usepackage{colortbl}
\usepackage{tabularx}
\usepackage{booktabs}
\usepackage{enumitem}
\usepackage{fancyhdr}
\usepackage{titlesec}
\usepackage{tocloft}
%\usepackage{microtype}  % Commented out due to font expansion issues

% --- PAQUETES DE MATEMÁTICAS Y FIGURAS ---
\usepackage{amsmath}
\usepackage{amssymb}
\usepackage{graphicx}
\usepackage{float}
\usepackage{caption}
\usepackage{subcaption}
\usepackage{longtable}

% --- DIAGRAMAS TIKZ (con protección Babel) ---
\usepackage{tikz}
\usetikzlibrary{arrows, arrows.meta, positioning, shapes.geometric, fit, backgrounds, calc, babel}

% --- COLORES INSTITUCIONALES UNAL Y TABLAS ---
\definecolor{unalnavy}{HTML}{003366}
\definecolor{unalblue}{HTML}{4682B4}
\definecolor{unalred}{HTML}{941B27}
\definecolor{categorybg}{HTML}{EBEBEB}     % Gris claro para encabezados
\definecolor{categorybgsub}{HTML}{DCDCDC}  % Gris medio para subcategorías
\definecolor{rowalt}{HTML}{F8F9FA}         % Gris ultra suave para alternar filas
\definecolor{phase1blue}{HTML}{E6F0FF}     % Fondo Fase 1 (Semestre)
\definecolor{phase2orange}{HTML}{FFF3E6}   % Fondo Fase 2 (Tesis)
\definecolor{bordergrey}{HTML}{B0B0B0}

\setlength{\arrayrulewidth}{0.8pt}

% --- FORMATO DE TÍTULOS EN NEGRO SOBRIO ---
\titleformat{\chapter}[hang]
  {\Huge\bfseries\color{black}}
  {\thechapter\hspace{15pt}}
  {0pt}
  {\Huge\bfseries}

\titleformat{\section}
  {\Large\bfseries\color{black}}
  {\thesection}{1em}{}

\titleformat{\subsection}
  {\large\bfseries\color{black}}
  {\thesubsection}{1em}{}

% --- ENCABEZADO Y PIE DE PÁGINA ---
% --- ENCABEZADO Y PIE DE PÁGINA (100% NEGRO) ---
\pagestyle{fancy}
\fancyhf{}
\lhead{\bfseries UNAL - DIEEC}
\rhead{\bfseries DR-VF-01 --- Cerca Virtual}
\rfoot{\small Página \thepage}
\lfoot{\small Septiembre, 2026 | Rev. 1.0}
\renewcommand{\headrulewidth}{0.8pt}
\renewcommand{\footrulewidth}{0.4pt}

% --- CONFIGURACIÓN DE HIPERVÍNCULOS (Cargar siempre al final) ---
\usepackage{hyperref}
\hypersetup{
    colorlinks=true,
    linkcolor=black,
    urlcolor=cyan,
    citecolor=unalnavy,
    pdftitle={Collar DR-VF-01 - Memoria Técnica Consolidada},
    pdfauthor={Proyecto DR-VF-01 - UNAL DIEEC}
}

% --- PERSONALIZACIÓN DE ÍNDICES ---
% Cargar tocloft sin opciones para evitar conflictos
\usepackage{tocloft}
% Configuración del Índice General
\renewcommand{\cfttoctitlefont}{\Huge\bfseries}
\renewcommand{\cftbeforetoctitleskip}{-1em}
\renewcommand{\cftaftertoctitleskip}{1em}

% Configuración del Índice de Figuras
\renewcommand{\cftloftitlefont}{\Huge\bfseries}
\renewcommand{\cftbeforeloftitleskip}{-1em}
\renewcommand{\cftafterloftitleskip}{1em}

% Configuración del Índice de Tablas
\renewcommand{\cftlottitlefont}{\Huge\bfseries}
\renewcommand{\cftbeforelottitleskip}{-1em}
\renewcommand{\cftafterlottitleskip}{1em}

% Espaciado entre entradas
\setlength{\cftbeforesecskip}{0.5em}
\setlength{\cftbeforesubsecskip}{0.3em}

% Asegurar que los índices se muestren correctamente
\setlength{\cftchapnumwidth}{2em}
\setlength{\cftsecnumwidth}{3em}
\setlength{\cftsubsecnumwidth}{4em}
\setlength{\cftfigpresnum}{\figurename~}
\setlength{\cfttabpresnum}{\tablename~}
\renewcommand{\cftchappresnum}{Capítulo~}
\setlength{\cftfigaftersnum}{:}
\setlength{\cfttabaftersnum}{:}